\chapter{Lezione 17 --- Esercizio su risposta al gradino, realizzazione minima e assegnamento dei poli}

\section{Esercizio: studio di una funzione di trasferimento discreta}

Consideriamo la funzione di trasferimento
\[
G(z)=\frac{z^2}{\left(z-\frac14\right)\left[\left(z-\frac12\right)^2+\frac14\right]}.
\]

Osserviamo subito che
\[
\left(z-\frac12\right)^2+\frac14
=
z^2-z+\frac12,
\]
quindi
\[
G(z)=\frac{z^2}{\left(z-\frac14\right)\left(z^2-z+\frac12\right)}.
\]

\subsection{Poli e modi del sistema}

I poli sono:
\[
z_1=\frac14,
\qquad
z_{2,3}=\frac12\pm j\,\frac12.
\]

Per i poli complessi coniugati conviene passare alla forma polare:
\[
r=\sqrt{\left(\frac12\right)^2+\left(\frac12\right)^2}
=\frac{\sqrt{2}}{2},
\qquad
\theta=\frac{\pi}{4}.
\]

Infatti
\[
\frac12\pm j\frac12
=
\frac{\sqrt2}{2}e^{\pm j\pi/4}.
\]

Quindi i modi naturali associati sono:
\[
\left(\frac14\right)^k,
\qquad
\left(\frac{\sqrt2}{2}\right)^k \cos\!\left(\frac{k\pi}{4}\right),
\qquad
\left(\frac{\sqrt2}{2}\right)^k \sin\!\left(\frac{k\pi}{4}\right).
\]

\section{Risposta al gradino}

Vogliamo calcolare la risposta al gradino unitario.

Nel dominio z:
\[
U(z)=\mathcal{Z}\{1(k)\}=\frac{z}{z-1}.
\]

Dunque
\[
Y(z)=G(z)U(z)
=
\frac{z^3}{(z-1)\left(z-\frac14\right)\left(z^2-z+\frac12\right)}.
\]

Per l'antitrasformata conviene scomporre in fratti semplici
\[
\frac{Y(z)}{z}
=
\frac{A}{z-1}
+
\frac{B}{z-\frac14}
+
C\,\frac{z-r\cos\theta}{z^2-2r\cos\theta\,z+r^2}
+
D\,\frac{r\sin\theta}{z^2-2r\cos\theta\,z+r^2}.
\]

Nel nostro caso
\[
r=\frac{\sqrt2}{2},
\qquad
\theta=\frac{\pi}{4},
\qquad
r\cos\theta=\frac12,
\qquad
r\sin\theta=\frac12,
\]
quindi
\[
\frac{Y(z)}{z}
=
\frac{A}{z-1}
+
\frac{B}{z-\frac14}
+
C\,\frac{z-\frac12}{z^2-z+\frac12}
+
D\,\frac{\frac12}{z^2-z+\frac12}.
\]

\subsection{Calcolo dei coefficienti}

Per i termini del primo ordine si può usare il teorema di Heaviside.

\paragraph{Coefficiente \(A\)}
\[
A=\lim_{z\to 1}\frac{Y(z)}{z}(z-1)
=
\frac{1}{\left(1-\frac14\right)\left(1-1+\frac12\right)}
=
\frac{1}{\frac34\cdot \frac12}
=
\frac{8}{3}.
\]

\paragraph{Coefficiente \(B\)}
\[
B=\lim_{z\to \frac14}\frac{Y(z)}{z}\left(z-\frac14\right)
=
\frac{1}{\left(\frac14-1\right)\left(\frac1{16}-\frac14+\frac12\right)}.
\]

Poiché
\[
\frac1{16}-\frac14+\frac12=\frac{5}{16},
\]
segue
\[
B
=
\frac{1}{-\frac34\cdot \frac5{16}}
=
-\frac{64}{15}
\cdot \frac{1}{16}
=
-\frac{4}{15}.
\]

\paragraph{Coefficienti \(C\) e \(D\)}

Scriviamo
\[
\frac{Y(z)}{z}
=
\frac{A}{z-1}
+
\frac{B}{z-\frac14}
+
C\,\frac{z-\frac12}{z^2-z+\frac12}
+
D\,\frac{\frac12}{z^2-z+\frac12}.
\]

Moltiplicando per il denominatore comune e uguagliando i coefficienti si ottiene:
\[
C=-\frac{12}{5},
\qquad
D=-\frac{4}{5}.
\]

\subsection{Antitrasformata z}

Usiamo le coppie note:
\[
\mathcal{Z}^{-1}\left\{\frac{z}{z-a}\right\}=a^k\,1(k),
\]
\[
\mathcal{Z}^{-1}\left\{
\frac{z(z-r\cos\theta)}{z^2-2r\cos\theta\,z+r^2}
\right\}
=
r^k\cos(k\theta)\,1(k),
\]
\[
\mathcal{Z}^{-1}\left\{
\frac{zr\sin\theta}{z^2-2r\cos\theta\,z+r^2}
\right\}
=
r^k\sin(k\theta)\,1(k).
\]

Moltiplicando di nuovo per \(z\), si ha
\[
Y(z)
=
\frac{8}{3}\frac{z}{z-1}
-
\frac{4}{15}\frac{z}{z-\frac14}
-
\frac{12}{5}\frac{z\left(z-\frac12\right)}{z^2-z+\frac12}
-
\frac{4}{5}\frac{z\cdot \frac12}{z^2-z+\frac12}.
\]

Quindi
\[
y(k)
=
\frac{8}{3}\,1(k)
-
\frac{4}{15}\left(\frac14\right)^k 1(k)
-
\frac{12}{5}\left(\frac{\sqrt2}{2}\right)^k
\cos\!\left(\frac{k\pi}{4}\right)1(k)
-
\frac{4}{5}\left(\frac{\sqrt2}{2}\right)^k
\sin\!\left(\frac{k\pi}{4}\right)1(k).
\]

Questa è la risposta al gradino.

\section{Realizzazione minima in equazioni di stato}

Vogliamo ora costruire una realizzazione minima del sistema, cioè del tipo
\[
\begin{cases}
x(k+1)=Ax(k)+Bu(k),\\
y(k)=Cx(k).
\end{cases}
\]

\subsection{Denominatore come polinomio}

Espandiamo il denominatore:
\[
\left(z-\frac14\right)\left(z^2-z+\frac12\right)
=
z^3-\frac54 z^2+\frac34 z-\frac18.
\]

Allora
\[
G(z)=\frac{z^2}{z^3-\frac54 z^2+\frac34 z-\frac18}.
\]

Poiché
\[
Y(z)=G(z)U(z),
\]
segue
\[
\left(z^3-\frac54 z^2+\frac34 z-\frac18\right)Y(z)=z^2U(z).
\]

\subsection{Equazione alle differenze}

Antitrasformando si può scrivere, ad esempio,
\[
y(k)-\frac54 y(k-1)+\frac34 y(k-2)-\frac18 y(k-3)=u(k-1).
\]

Questa equazione suggerisce una realizzazione canonica di controllo di ordine \(3\).

\subsection{Forma canonica di controllo}

Poiché il denominatore ha grado \(3\), la dimensione minima dello stato è \(3\).  
Una realizzazione minima è quindi:
\[
x(k+1)=
\begin{bmatrix}
0 & 1 & 0\\
0 & 0 & 1\\
\frac18 & -\frac34 & \frac54
\end{bmatrix}x(k)
+
\begin{bmatrix}
0\\
0\\
1
\end{bmatrix}u(k),
\]
\[
y(k)=
\begin{bmatrix}
0 & 0 & 1
\end{bmatrix}x(k).
\]

\section{Osservazione sulle realizzazioni minime}

Dato un sistema in forma di stato, esiste un'unica funzione di trasferimento associata.

Viceversa, data una funzione di trasferimento, esistono in generale infiniti sistemi in forma di stato che la realizzano.

Tra tutte queste realizzazioni, si chiamano \textbf{realizzazioni minime} quelle con dimensione dello stato minima possibile.

Una realizzazione minima è, per costruzione, completamente:
\begin{itemize}
    \item raggiungibile (o controllabile);
    \item osservabile.
\end{itemize}

In particolare, i poli della funzione di trasferimento corrispondono agli autovalori raggiungibili e osservabili del sistema.

Quindi, una volta ottenuta una realizzazione minima, sono soddisfatte le ipotesi necessarie
per applicare le tecniche di assegnamento dei poli mediante retroazione dello stato.

\section{Progetto in ciclo chiuso: scelta dei poli desiderati}

Il sistema ha un polo reale in
\[
\frac14
\]
e una coppia di poli complessi coniugati in
\[
\frac12\pm j\frac12.
\]

Non essendoci specifiche particolari, scegliamo poli desiderati:
\begin{itemize}
    \item reali;
    \item più vicini a zero;
    \item quindi più veloci dei poli originari.
\end{itemize}

Una scelta semplice è:
\[
z_1^\star=\frac15,
\qquad
z_2^\star=\frac1{10},
\qquad
z_3^\star=\frac1{20}.
\]

\subsection{Polinomio caratteristico desiderato}

Il polinomio caratteristico desiderato è
\[
p_c^\star(\lambda)
=
\left(\lambda-\frac15\right)
\left(\lambda-\frac1{10}\right)
\left(\lambda-\frac1{20}\right).
\]

Sviluppando:
\[
p_c^\star(\lambda)
=
\lambda^3
-\frac{7}{20}\lambda^2
+\frac{7}{200}\lambda
-\frac{1}{1000}.
\]

\section{Retroazione completa dello stato}

Introduciamo una retroazione completa dello stato del tipo
\[
u(k)=-H\,x(k),
\qquad
H=
\begin{bmatrix}
h_1 & h_2 & h_3
\end{bmatrix}.
\]

La dinamica in ciclo chiuso diventa
\[
x(k+1)=(A-BH)x(k).
\]

Poiché
\[
B=
\begin{bmatrix}
0\\0\\1
\end{bmatrix},
\]
si ha
\[
BH=
\begin{bmatrix}
0\\0\\1
\end{bmatrix}
\begin{bmatrix}
h_1 & h_2 & h_3
\end{bmatrix}
=
\begin{bmatrix}
0 & 0 & 0\\
0 & 0 & 0\\
h_1 & h_2 & h_3
\end{bmatrix}.
\]

Quindi
\[
A-BH=
\begin{bmatrix}
0 & 1 & 0\\
0 & 0 & 1\\
\frac18-h_1 & -\frac34-h_2 & \frac54-h_3
\end{bmatrix}.
\]

Essendo ancora in forma canonica di controllo, il suo polinomio caratteristico è
\[
\lambda^3
+\left(-\frac54+h_3\right)\lambda^2
+\left(\frac34+h_2\right)\lambda
+\left(-\frac18+h_1\right).
\]

Imponiamo l'uguaglianza con il polinomio desiderato:
\[
\lambda^3
+\left(-\frac54+h_3\right)\lambda^2
+\left(\frac34+h_2\right)\lambda
+\left(-\frac18+h_1\right)
=
\lambda^3-\frac{7}{20}\lambda^2+\frac{7}{200}\lambda-\frac{1}{1000}.
\]

Da cui:
\[
-\frac54+h_3=-\frac{7}{20},
\qquad
\frac34+h_2=\frac{7}{200},
\qquad
-\frac18+h_1=-\frac{1}{1000}.
\]

Si ottiene quindi
\[
h_3=\frac{9}{10},
\qquad
h_2=-\frac{143}{200},
\qquad
h_1=\frac{31}{250}.
\]

Dunque la matrice di guadagno è
\[
H=
\begin{bmatrix}
\frac{31}{250} & -\frac{143}{200} & \frac{9}{10}
\end{bmatrix}.
\]

\section{Commento finale}

Questa lezione mette insieme tre aspetti importanti.

\begin{enumerate}
    \item Data una funzione di trasferimento discreta, si possono ricavare:
    \begin{itemize}
        \item poli;
        \item modi naturali;
        \item risposta al gradino;
        \item una realizzazione minima in forma di stato.
    \end{itemize}

    \item Una realizzazione minima è la più conveniente per il progetto, perché è:
    \begin{itemize}
        \item completamente raggiungibile;
        \item completamente osservabile.
    \end{itemize}

    \item Se la realizzazione è in forma canonica di controllo, l'assegnamento dei poli in ciclo chiuso
    è particolarmente semplice: basta confrontare i coefficienti del polinomio caratteristico.
\end{enumerate}

\section{Sintesi della lezione}

I risultati principali sono:
\begin{itemize}
    \item i poli di
    \[
    G(z)=\frac{z^2}{\left(z-\frac14\right)\left(z^2-z+\frac12\right)}
    \]
    sono
    \[
    \frac14,\qquad \frac12\pm j\frac12;
    \]

    \item la risposta al gradino è
    \[
    y(k)
    =
    \frac{8}{3}\,1(k)
    -
    \frac{4}{15}\left(\frac14\right)^k 1(k)
    -
    \frac{12}{5}\left(\frac{\sqrt2}{2}\right)^k
    \cos\!\left(\frac{k\pi}{4}\right)1(k)
    -
    \frac{4}{5}\left(\frac{\sqrt2}{2}\right)^k
    \sin\!\left(\frac{k\pi}{4}\right)1(k);
    \]

    \item una realizzazione minima è
    \[
    x(k+1)=
    \begin{bmatrix}
    0 & 1 & 0\\
    0 & 0 & 1\\
    \frac18 & -\frac34 & \frac54
    \end{bmatrix}x(k)
    +
    \begin{bmatrix}
    0\\0\\1
    \end{bmatrix}u(k),
    \qquad
    y(k)=
    \begin{bmatrix}
    0 & 0 & 1
    \end{bmatrix}x(k);
    \]

    \item scegliendo i poli desiderati
    \[
    \frac15,\quad \frac1{10},\quad \frac1{20},
    \]
    il polinomio caratteristico desiderato è
    \[
    \lambda^3-\frac{7}{20}\lambda^2+\frac{7}{200}\lambda-\frac{1}{1000};
    \]

    \item la retroazione completa dello stato
    \[
    u(k)=-H x(k)
    \]
    con
    \[
    H=
    \begin{bmatrix}
    \frac{31}{250} & -\frac{143}{200} & \frac{9}{10}
    \end{bmatrix}
    \]
    assegna esattamente tali poli al sistema in ciclo chiuso.
\end{itemize}